% Part: computability
% Chapter: computability-theory
% Section: no-universal-function

\documentclass[../../../include/open-logic-section]{subfiles}

\begin{document}

\olfileid{cmp}{thy}{nou}
\olsection{Ora Ana Fungsi Komputabel Universal}

% OLPL-135: Sumber beku nyebut fungsi parsial sadurunge total kanggo fungsi parsial, nanging teorema sadurunge lan panjelasan ing ngisor iki nyebut universal lan ora total.
Sanajan ana fungsi komputabel parsial kang
universal kanggo fungsi komputabel parsial,
ora ana fungsi komputabel total kang universal
kanggo fungsi komputabel total.

\begin{thm}
\ollabel{thm:no-univ}
Ora ana fungsi komputabel universal. Kanthi
tembung liya, fungsi $\fn{Un}'(k, x)$ sembarang
kang nduwé sipat iki: yen $f(x)$ fungsi komputabel
total, ana wilangan asli~$k$ kang ndadekake
$f(x) = \fn{Un}'(k,x)$ kanggo saben~$x$,
fungsi kasebut ora komputabel.
\end{thm}

\begin{proof}
Buktine yaiku argumen diagonal kang prasaja:
yen $\fn{Un}'(k,x)$ total lan komputabel,
mula
\[
d(x) = \fn{Un}'(x, x) + 1
\]
uga total lan komputabel. Nanging miturut
definisi, $d(k)$ ora padha karo
$\fn{Un}'(k,k)$. Mula, kanggo saben $k$,
nilai $d(x)$ lan~$\fn{Un}'(k, x)$ beda
ing saora-orane siji~$x$, yaiku $x = k$.
\end{proof}

\begin{explain}
\olref[uni]{thm:univ-comp} ing ndhuwur nuduhake
yen argumen diagonal iki bisa diendhani,
nanging mung kanthi ngidini fungsi universal
iku parsial. Tegese, $\fn{Un}$ universal
kanggo fungsi komputabel total, nanging
fungsi kasebut dhewe ora total. Argumen
diagonal ora lumaku ing kasus parsial.
\end{explain}

\begin{prob}
  Kanggo mangerteni ngapa argumen diagonal
  ing bukti \olref{thm:no-univ} ora lumaku
  ing kasus parsial, gatekna fungsi
  $f(x) \simeq \fn{Un}(x,x)+1$. Apa fungsi
  iki komputabel parsial? Yen mangkono,
  fungsi iki nduwé indeks~$e$, yaiku
  $f(x) \simeq \fn{Un}(e,x)$. Apa kang
  bisa kokandharake bab~$f(e)$?
\end{prob}

\end{document}
